\documentclass[a4paper]{report}
% Shared preamble for the English probability and statistics documents.
% Based on the original course preamble. Compile with pdfLaTeX.
\usepackage[top=1.5cm,bottom=2cm,left=2cm,right=2cm]{geometry}
\usepackage[T1]{fontenc}
\usepackage{lmodern}
\usepackage[utf8]{inputenc}
\usepackage[english]{babel}

\usepackage{amsmath}
\usepackage{amssymb}
\usepackage{amsthm}
\usepackage{mathrsfs}
\usepackage{mathtools}
\usepackage{etoolbox}
\usepackage{tikz}
\usetikzlibrary{shapes.multipart,arrows.meta}
\usepackage{xcolor}
\usepackage{enumerate}
\usepackage[shortlabels]{enumitem}
\usepackage{hyperref}
\usepackage{version}
\usepackage{bbold}
\usepackage{fontawesome5}
\usepackage{fancyhdr}
\lfoot{}
\usepackage{natbib}
\usepackage{graphicx}
\usepackage{pdfpages}
\usepackage{authblk}
    
\usepackage{blindtext}
\usepackage{multicol}
\usepackage{collcell}
\usepackage{wrapfig}
\usepackage{blkarray}
\usepackage{cancel}
\usepackage{appendix}
\usepackage{float}
\makeatletter
\input{tex/assets/slashbox.sty}
\makeatother
\hypersetup{unicode,colorlinks,citecolor=red,filecolor=black,
  linkcolor=red,urlcolor=cyan}
\setlength{\columnsep}{1cm}
\pagenumbering{gobble}

\def\numtext#1{%
  \ifcase#1\or one\or two\or three\or four\or five\or six\or seven\or eight\or nine\or ten\or
  eleven\or twelve\or thirteen\or fourteen\or fifteen\or sixteen\or seventeen\or eighteen\or nineteen\or twenty\or Lots\fi}

\newcommand{\Byte}[1]{%
  \def\nodecontents{}%
  \foreach[count=\l] \k in {#1} {\xappto{\nodecontents}{\noexpand\nodepart{\numtext{\l}}\k}}%
  \node[rectangle split,rectangle split parts=8,rectangle split horizontal,draw] at (2,2){\nodecontents};}

\newcommand{\D}{{\mathrm{d}}}
\newcommand{\cE}{{\mathcal{E}}}
\newcommand{\cG}{{\mathcal{G}}}
\newcommand{\cX}{{\mathcal{X}}}
\newcommand{\cY}{{\mathcal{Y}}}
\newcommand{\cB}{{\mathcal{B}}}
\newcommand{\cF}{{\mathcal{F}}}
\newcommand{\cO}{{\mathcal{O}}}
\newcommand{\cS}{{\mathcal{S}}}
\newcommand{\cT}{{\mathcal{T}}}
\newcommand{\given}{\, \big| \,}
\newcommand{\cL}{{\mathscr{L}}}
\newcommand{\cD}{{\mathcal{D}}}
\newcommand{\cA}{{\mathscr{A}}}
\newcommand{\cP}{{\mathscr{P}}}
\newcommand{\cR}{{\mathscr{R}}}
\newcommand{\A}{{\mathcal A}}
\newcommand{\one}{\mathbb{1}}
\newcommand{\mathbbm}[1]{\mathbb{#1}}

\providecommand{\C}{}
\renewcommand{\C}{\mathbb C}
\newcommand{\R}{\mathbb R}
\newcommand{\Q}{\mathbb Q}
\newcommand{\N}{\mathbb N}
\newcommand{\T}{\mathbb T}
\newcommand{\F}{\mathbb F}
\renewcommand{\P}{\mathbb P}
\newcommand{\Z}{\mathbb Z}
\newcommand{\E}{\mathbb {E}}
\newcommand{\Pbb}[1]{\P\left( #1\right)}
\newcommand{\Ebb}[1]{\E\left[ #1\right]}

\newcommand{\lb}{\left (}
\newcommand{\rb}{\right )}
\newcommand{\lbb}{\left [}
\newcommand{\rbb}{\right ]}
\newcommand{\labs}{\left |}
\newcommand{\rabs}{\right |}
\newcommand{\lbrb}[1]{\lb #1 \rb}
\newcommand{\lbbrbb}[1]{\lbb#1\rbb}
\newcommand{\labsrabs}[1]{\labs#1\rabs}
\newcommand{\lbbrb}[1]{\lbb#1\rb}
\newcommand{\lbrbb}[1]{\lb#1\rbb}
\newcommand{\langlerangle}[1]{\langle#1\rangle}
\newcommand{\lbcurly}{\{}
\newcommand{\rbcurly}{\}}
\newcommand{\curly}[1]{\lbcurly#1\rbcurly}
\newcommand{\abs}{\labsrabs}
\newcommand{\ind}[1]{1_{#1}}
\newcommand{\indep}{\perp \!\!\! \perp}
\renewcommand{\d}{\, \mathrm{d}}

\newtheorem{theorem}{Theorem}
\newtheorem{problemCnt}{Lemma}
\theoremstyle{definition}
\newtheorem{problem}[problemCnt]{Problem}
\newtheorem{problembook}[problemCnt]{\faBook \,  Problem}
\newtheorem{problempencil}[problemCnt]{\faPencil* \,  Problem}
\newtheorem{problemstar}[problemCnt]{\faStar\,  Problem}
\newtheorem{problemstarpencil}[problemCnt]{\faStar \, \faPencil* \,   Problem}
\newtheorem{probleminteresting}[problemCnt]{\faSmile[regular] \,   Problem}
\newtheorem{solution}{Solution}

\newcommand{\fwcell}[1]{\makebox[\arraystretch\normalbaselineskip]{$#1$}}
\newenvironment{enumeratei}{\begin{enumerate}[(i)]}{\end{enumerate}}
\newenvironment{enumeratea}{\begin{enumerate}[(a)]}{\end{enumerate}}
\newenvironment{enumeraten}{\begin{enumerate}[\upshape 1.]}{\end{enumerate}}

% Compatibility helpers needed by a few files in the combined archive.
\AtBeginDocument{\noindent{\footnotesize English translation prepared with ChatGPT (OpenAI, GPT-6), 26 October 2026, for translation of the Bulgarian text into English and LaTeX formatting.
The Bulgarian originals remain the reference; translation may contain errors.}\par\smallskip}
\newcommand{\MissingGraphic}[1]{%
  \begin{center}\fbox{\parbox{0.82\linewidth}{\centering Missing source image: \texttt{#1}}}\end{center}}



\begin{document}
\noindent Probability Theory and Mathematical Statistics, 2023/2024 \hfill Homework 3
\hrulefill
\begin{it}The maximum total across all homework assignments this semester is 50 points, corresponding to a bonus of 0.5 on the tutorial grade. Unless otherwise stated, $m,n,k\in\mathbb N=\{1,2,3,\ldots\}$; define $\mathbb N_0=\mathbb N\cup\{0\}$.\end{it}
\begin{problem} (2 points)
A person starts at $n\in\mathbb N$ on the number line, moving one step right with probability $p>1/2$ and one step left with probability $1-p$. Let $p_n$ be the probability of reaching 0 when starting at $n$. Express $p_1,p_2,p_3$ in terms of $p$.
\end{problem}

\begin{problem} (2 points)
A rabbit starts at 0 and makes independent jumps uniformly distributed on $[0,1]$ in the positive direction. The interval $[1-x,1]$ is a trap of length $x\in[0,1]$. What is the probability that the rabbit jumps over the trap?
\end{problem}
\begin{problem} (6 points; simulating random variables)
Assume that we can readily simulate $U(0,1)$ variables. This may be a mathematical task, constructing deterministic sequences that ``look like'' independent uniform samples (pseudorandom numbers), or a technological task using physical phenomena.

We study three classical algorithms for simulating a value $X$ from a specified distribution: inversion of $F_X$, rejection sampling, and the Box--Muller method.

\textbf{Method 1}
\begin{enumerate}
\item Let $U\sim U(0,1)$. What is the distribution of $F_X^{-1}(U)$?\footnote{If $F_X$ is not strictly increasing, its ordinary inverse may not exist everywhere. One usually uses the generalized inverse $F_X^{-1}(t)=\inf\{x\in\mathbb R:F_X(x)\geq t\}$. For our purposes you may assume an ordinary inverse.} Use this to propose simulations of:
\begin{enumerate}
\item $X\sim\mathrm{Exp}(\lambda)$, $\lambda>0$;
\item $X\sim\mathrm{Cauchy}(c)$, $c>0$, with $f_X(x)=c/[\pi(c^2+x^2)]$;
\item $X\sim\mathrm{Pareto}(\theta)$, $\theta>0$, with $f_X(x)=\theta x^{-1-\theta}\mathbbm1_{\{x\geq1\}}$.
\end{enumerate}
\item What phenomena do these distributions model, and why? Generate samples and illustrate the law of large numbers and central limit theorem in all three cases.
\item What difficulties arise if this method is used for normal or gamma distributions?
\end{enumerate}
\textbf{Method 2}
The following algorithm, associated with ideas of von Neumann from 1949--1950, simulates a variable with density $f_X$:
\begin{itemize}
\item Choose a variable $Y$ that is easy to simulate, for example by Method 1, with density $f_Y$ satisfying $f_X(x)/f_Y(x)\leq c$ for all $x$ and some $c>0$.
\item Generate a value of $Y$.
\item Independently generate $U\sim U(0,1)$.
\item If $U\leq f_X(Y)/[cf_Y(Y)]$, accept $Y$ as a sample from $X$; otherwise reject it and repeat.
\end{itemize}
\begin{enumerate}
\item Prove that the algorithm is correct.
\item Generate $X\sim N(\mu,\sigma^2)$. \emph{One possibility is to use an exponential proposal after a suitable transformation.}
\item Let $X\sim\Gamma(\alpha,1)$, with density $f_X(x)=x^{\alpha-1}e^{-x}\mathbbm1_{\{x>0\}}/\Gamma(\alpha)$.
\begin{enumerate}
\item For $\alpha<1$, consider $Y$ with density
\[f_Y(x)=\frac{\alpha e}{\alpha+e}\left(x^{\alpha-1}\mathbbm1_{\{0<x<1\}}+e^{-x}\mathbbm1_{\{x\geq1\}}\right).\]
\item Show that this is a density. How would you simulate $Y$ using Method 1?
\item Can Method 2 be used to generate $\Gamma(\alpha,1)$ variables for $\alpha<1$?
\item How would you simulate $\Gamma(\alpha,\beta)$ using this method and properties of the gamma distribution?
\end{enumerate}
\end{enumerate}
\textbf{Method 3}
\begin{enumerate}
\item Let $X^2\sim\mathrm{Exp}(1/2)$ and $\Theta\sim U(0,2\pi)$ be independent. Find the distributions of $Z_1=R\cos\Theta$ and $Z_2=R\sin\Theta$. Are they independent?
\item How would you use this to simulate $n$ independent $N(\mu,\sigma)$ variables?
\end{enumerate}
A frequently used distribution (why?) is the lognormal: $X\sim\mathrm{Lognormal}(\mu,\sigma^2)$ when $\ln X\sim N(\mu,\sigma)$. Propose three ways to simulate it.
\end{problem}
\begin{problem} (2 points)
Let $X$ take values in $\mathbb N_0=\{0,1,2,\ldots\}$, with $\E[X]=1$, $\E[X^2]=2$, $\E[X^3]=5$. Find the smallest possible value of $\P(X=0)$.
\end{problem}
\begin{problem} (3 points)
Consider discrete real-valued random variables with finite support, that is, with values bounded above and below by constants. Write $X\leq_cY$ if $\E[f(X)]\leq\E[f(Y)]$ for every convex $f:\mathbb R\to\mathbb R$.
\begin{enumerate}
\item Prove $f(\E[X])\leq\E[f(X)]$ for such $X$ and convex $f$.
\item Give an example with $X\leq_cY$ but $X\neq Y$.
\item Prove that $X\leq_cY$ implies $\E[X]=\E[Y]$ and $D(X)\leq D(Y)$.
\item Prove that $X\leq_cY$ holds if and only if $\E[X]=\E[Y]$ and, for every $a\in\mathbb R$,
\[\int_a^\infty\P(X\geq x)\,\mathrm dx\leq\int_a^\infty\P(Y\geq x)\,\mathrm dx.\]
\end{enumerate}
\end{problem}
\begin{problem} (3 points)
Let $E$ be finite, $V:E\to2^E$, and $f:E\to\mathbb R$. Call $a\in E$ a local minimum if $f(a)\leq f(b)$ for every $b\in V(a)$.
Let $M\geq\sqrt{|E|}$ be an integer and $b_1,\ldots,b_M$ independent and uniform on $E$. Choose $k$ so $f(b_k)=\min_{1\leq i\leq M}f(b_i)$.
Let $u_0=b_k$ and, for every $n\geq0$, set $u_{n+1}=u_n$ if $u_n$ is a local minimum; otherwise choose $u_{n+1}\in V(u_n)$ with $f(u_{n+1})<f(u_n)$. Prove that $u_M$ is a local minimum with probability at least $1/2$.
\end{problem}
\begin{problem} (3 points)
Let $d,n\in\mathbb N$, $n\geq3$, $G=(\mathbb Z/n\mathbb Z)^d$ and $S=\{\pm e_i:1\leq i\leq d\}$, where $e_i$ has coordinate $\overline1$ in position $i$ and zeros elsewhere. Let $f:G\to\mathbb R$ and $X$ be uniform on $G$. Prove
\[\E[|f(X)-\E[f(X)]|]\leq\frac{nd}{2}\max_{s\in S}\E[|f(X)-f(X+s)|].\]
\end{problem}
\begin{problem} (3 points)
Let $X_n$ be iid with values in $\mathbb N$. Assume $\P(X_1=0)\P(X_1=1)\neq0$ and set $S_n=X_1+\cdots+X_n$. Prove $\P(4\text{ divides }S_n)\to1/4$ as $n\to\infty$.
\end{problem}
\begin{problem} (3 points)
Let $Z\sim N(0,1)$. Prove for $x>0$ that
\[\frac{\sqrt{x^2+4}-x}{2e^{x^2/2}}<\sqrt{2\pi}\P(Z>x)<\frac1{xe^{x^2/2}}.\]
Find the first four terms of the asymptotic expansion of $f(x)=e^{x^2/2}\P(Z>x)$ as $x\to\infty$.
\end{problem}
\begin{problem} (2 points)
Let $n>1$ and let $A,B,C$ be independent and uniform on $\{0,1\}^n$. For $n>2$, calculate $p_n=\P(ABC\text{ is an equilateral triangle})$ and determine its asymptotic behaviour.
\end{problem}

\end{document}
